\section{Framework: Policy Evaluation When Overlap Fails by Design}\label{sec:framework}

\subsection{Treatment, menus and potential outcomes}

Index scheduled FOMC meetings by $t$. The policy decision is $D_t\in\cD=\{-1,0,1\}$ (ease, leave unchanged,
tighten the federal funds rate target). The Bluebook menu is $X_t\in\cX$, a non-empty subset of $\cD$; in the data
$\cX=\{\{-1\},\{0\},\{1\},\{-1,0\},\{0,1\},\{-1,0,1\}\}$. Other information available to the Committee (lagged
macroeconomic data, market expectations, staff forecasts) is collected in $Z_t\in\mathbb{R}^k$, and
$W_t=(X_t,Z_t)$.

For an outcome series $y$ (log industrial production, the unemployment rate, the log PCE price index), let
$Y_{t,h}=y_{t+h}-y_t$ be its cumulative change over the $h$ months following the meeting. Each decision $d$
induces a potential outcome $Y_{t,h}(d)$, and $Y_{t,h}=\sum_{d\in\cD}\1\{D_t=d\}\,Y_{t,h}(d)$. The parameter of
interest is the average dynamic effect of an increase ($d=1$) or a decrease ($d=-1$) relative to leaving the target
unchanged,
\begin{equation}\label{eq:theta}
  \theta_d(h)=\E\bigl[Y_{t,h}(d)-Y_{t,h}(0)\bigr],\qquad d\in\{-1,1\},\ h=1,\dots,24 .
\end{equation}
Expectations are taken over the stationary distribution of meetings. I drop the $t$ and $h$ subscripts when they
are not needed. The propensity score is $p_d(w)=\Prob(D=d\mid W=w)$.

\subsection{Assumptions}

\begin{assumption}[Selection on observables]\label{as:unconf}
$Y_{h}(d)\indep D\mid W$ for all $d\in\cD$ and all $h$.
\end{assumption}

\begin{assumption}[Menu overlap]\label{as:menu}
For all $w=(x,z)$ in the support of $W$: $p_d(x,z)>0$ if and only if $d\in x$.
\end{assumption}

\begin{assumption}[Bounded outcomes]\label{as:bounded}
$Y_{h}(d)\in[K_{1h},K_{2h}]$ almost surely, for all $d$ and $h$, with known $K_{1h}<K_{2h}$.
\end{assumption}

Assumption~\ref{as:unconf} is the identifying assumption of \citet{angrist2018}. It says that once we condition
on what the Committee knew, including what it was offered, the remaining variation in its decision is as good as
random with respect to the economy's future path. Assumption~\ref{as:menu} replaces the usual overlap condition,
$p_d(w)>0$ for all $d$ and $w$. It allows the propensity score to be zero, but only in the way the Bluebook
dictates: the Committee does not take an action that is not on the menu (which holds without exception in the
data, Table~\ref{tab:menus}), and any action on the menu has some chance of being taken. The ``only if'' part is
testable and holds. The ``if'' part is an assumption. Assumption~\ref{as:bounded} is a support restriction of
the kind used by \citet{manski1990}. It bites only on the cells in which the data are silent.

\subsection{What the data identify}

Write $\pi_x=\Prob(X=x)$ and $\mu_d(x)=\E[Y(d)\mid X=x]$. Conditioning on the menu,
\begin{equation}\label{eq:decomp}
  \theta_d=\sum_{x\in\cX}\pi_x\,\bigl(\mu_d(x)-\mu_0(x)\bigr).
\end{equation}
Under Assumptions~\ref{as:unconf}--\ref{as:menu}, for any $a\in x$ the conditional mean $\mu_a(x)$ is identified by
inverse probability weighting within the menu:
\begin{equation}\label{eq:ipw}
  \mu_a(x)=\E\!\left[\frac{\1\{D=a\}\,Y}{p_a(X,Z)}\ \middle|\ X=x\right],\qquad a\in x,
\end{equation}
because $p_a(x,Z)>0$ for every $Z$ when $a\in x$. When $a\notin x$, the data never show $Y(a)$ in menu $x$ and
$\mu_a(x)$ is not identified. For a given $d$, partition the menus into four groups:
\begin{align*}
\cA_d&=\{x:\,d\in x,\,0\in x\}, & \cB_d&=\{x:\,0\in x,\,d\notin x\},\\
\cC_d&=\{x:\,d\in x,\,0\notin x\}, & \cN_d&=\{x:\,d\notin x,\,0\notin x\}.
\end{align*}
For $d=1$ these are $\cA_1=\{\{0,1\},\{-1,0,1\}\}$, $\cB_1=\{\{0\},\{-1,0\}\}$, $\cC_1=\{\{1\}\}$ and
$\cN_1=\{\{-1\}\}$. The overlap region $\cA_d$ is where the effect is point-identified. In $\cB_d$ the status-quo mean
is identified but the treated mean is not; in $\cC_d$ the reverse; in $\cN_d$ neither.

\begin{proposition}[Sharp identified set]\label{prop:bounds}
Under Assumptions~\ref{as:unconf}--\ref{as:bounded}, the identified set for $\theta_d(h)$ is the interval
$[L_d,U_d]$ with
\begin{align*}
L_d&=\sum_{x\in\cA_d}\pi_x\bigl(\mu_d(x)-\mu_0(x)\bigr)+\sum_{x\in\cB_d}\pi_x\bigl(K_1-\mu_0(x)\bigr)\\
&\qquad+\sum_{x\in\cC_d}\pi_x\bigl(\mu_d(x)-K_2\bigr)+\sum_{x\in\cN_d}\pi_x\,(K_1-K_2),\\
U_d&=\sum_{x\in\cA_d}\pi_x\bigl(\mu_d(x)-\mu_0(x)\bigr)+\sum_{x\in\cB_d}\pi_x\bigl(K_2-\mu_0(x)\bigr)\\
&\qquad+\sum_{x\in\cC_d}\pi_x\bigl(\mu_d(x)-K_1\bigr)+\sum_{x\in\cN_d}\pi_x\,(K_2-K_1).
\end{align*}
Its width is $U_d-L_d=(K_2-K_1)\bigl[\pi(\cB_d)+\pi(\cC_d)+2\,\pi(\cN_d)\bigr]$.
\end{proposition}

The proof is in Appendix~\ref{app:proofs}. The width is always below the width of the no-information bound,
$2(K_2-K_1)$, so the set is informative whenever some menu contains both $d$ and $0$ or one of them. It is also
narrower than the naive bound that treats every cell outside $\cA_d$ as uninformative, whose width is
$2(K_2-K_1)\bigl[1-\pi(\cA_d)\bigr]$, because the identified half of each one-sided cell is used. The 2020
version of this thesis used the naive width.

\begin{remark}[Additional covariates]
Conditioning on $Z$ in \eqref{eq:ipw} is needed for Assumption~\ref{as:unconf}, but the bounds only require the
menu-level means. If $X$ happens to contain all the information in $Z$ that predicts $D$, the propensity score
does not depend on $Z$ and \eqref{eq:ipw} reduces to a difference in means within menu; nothing in
Proposition~\ref{prop:bounds} changes (Appendix~\ref{app:proofs}).
\end{remark}

\subsection{A sensitivity parameter for the missing menus}

The worst case in Proposition~\ref{prop:bounds} lets the average effect in menus outside the overlap region be as
extreme as the outcome's support permits. A more useful statement bounds how different those effects can be from
the identified effect in the overlap region,
\[
\theta_d^{\cA}=\E\bigl[Y(d)-Y(0)\mid X\in\cA_d\bigr]=\frac{\sum_{x\in\cA_d}\pi_x\bigl(\mu_d(x)-\mu_0(x)\bigr)}{\pi(\cA_d)}.
\]

\begin{assumption}[$\lambda$-bounded heterogeneity]\label{as:lambda}
For some $\lambda\ge0$ and all $x\notin\cA_d$: $\bigl|\E[Y(d)-Y(0)\mid X=x]-\theta_d^{\cA}\bigr|\le\lambda$.
\end{assumption}

\begin{proposition}\label{prop:lambda}
Under Assumptions~\ref{as:unconf}--\ref{as:lambda} the identified set is obtained from Proposition~\ref{prop:bounds}
by intersecting, for each $x\notin\cA_d$, the cell-level interval with $[\theta^{\cA}_d-\lambda,\theta^{\cA}_d+\lambda]$.
When the support constraint does not bind this is
\[
\Theta_d(\lambda)=\Bigl[\theta_d^{\cA}-\lambda\bigl(1-\pi(\cA_d)\bigr),\ \theta_d^{\cA}+\lambda\bigl(1-\pi(\cA_d)\bigr)\Bigr],
\]
and the sign of $\theta_d$ is identified if and only if $\lambda<\lambda^\ast_d=|\theta_d^{\cA}|/\bigl(1-\pi(\cA_d)\bigr)$.
\end{proposition}

At $\lambda=0$ effects are homogeneous across menus and $\theta_d=\theta_d^{\cA}$; as $\lambda$ grows the set
widens until it reaches the worst case. The \emph{breakdown value} $\lambda^\ast_d$ is the smallest degree of
heterogeneity across menus that would overturn the sign of the estimated effect, measured in the units of the
outcome. It gives a transparent way to judge how much the conclusion depends on extrapolating from the overlap
region.
